\documentclass[10pt,a4paper]{amsart}
\usepackage[margin=19mm]{geometry}
\usepackage{amsmath,amssymb,mathtools}
\usepackage[scr=rsfs,cal=euler]{mathalfa}
\usepackage{xfrac}
\usepackage[unicode,hidelinks,bookmarks=false]{hyperref}
\newcommand{\ie}{i.e.\ }
\newcommand{\cf}{cf.\ }
\newcommand{\resp}{respectively\ }
\newcommand{\Subsection}{\subsection}
\providecommand{\nobreakdash}{\nobreak\hbox{-}}
\setlength{\emergencystretch}{2em}
\multlinegap0pt
\usepackage{xcolor}
\usepackage{algorithm}\usepackage{algpseudocode}
\usepackage{tikz}
\usepackage{siunitx}
\usepackage{mathtools}
\usepackage{booktabs}
\newtheorem{proposition}{Proposition}
\newcommand{\D}{\mathbf D}
\newcommand{\Dhat}{\widehat{\D}}
\newcommand{\Hhat}{\widehat{\mathbf H}}
\newcommand{\V}{\mathbf V}
\newcommand{\bPhi}{\mathbf{\Phi}}
\newcommand{\dd}{\rm d}
\newcommand{\p}{\mathbf p}
\newcommand{\phat}{\widehat{\p}}
\newcommand{\phatdot}{\mathbf{\dot{\widehat p}}}
\newcommand{\q}{\mathbf q}
\newcommand{\qhat}{\widehat{\q}}
\newcommand{\qhatdot}{\mathbf{\dot{\widehat q}}}
\newcommand{\real}{\mathbb{R}}
\newcommand{\w}{\mathbf w}
\newcommand{\what}{\widehat{\w}}
\title{Structure-preserving Lift \& Learn: Scientific machine learning for nonlinear conservative partial differential equations}
\author{Harsh Sharma, Juan Diego Draxl Giannoni, Boris Kramer}
\date{Source: arXiv:2507.00301v2 (CC BY 4.0)}
\begin{document}
\maketitle

\noindent\emph{Source and licence.} Structure-preserving Lift \& Learn: Scientific machine learning for nonlinear conservative partial differential equations. Version arXiv:2507.00301v2 (CC BY 4.0), distributed by arXiv under the Creative Commons Attribution 4.0 International licence, http://creativecommons.org/licenses/by/4.0/, as recorded in the arXiv licence metadata for this exact version and shown on the abstract page https://arxiv.org/abs/2507.00301v2. Full licence notice: \texttt{licenses/arxiv-2507.00301-CC-BY-4.0.txt}.\par\smallskip

\noindent\textit{Editorial scope.} Counted unit: the article's proposition eq:pertrubed\_energy (source0.tex lines 1341--1347). The article's complete original proof of it is delivered with it (source0.tex lines 1348--1366, one \texttt{proof} environment of 19 lines). No mathematics is written, repaired or re-derived: the statement and the proof are the article's own lines, in the article's own notation, and no retained line is substituted or re-flowed.  Retained uncounted as the source-local context the proof needs: lines 1266--1272; lines 1283--1286; lines 1322--1328.  Nothing was omitted from any retained span, so the package carries no omission record.  No retained line contains a citation, so the wrapper carries no bibliography and no bibliography entry.  No retained line contains a figure environment, an image-inclusion command or drawing code, so no image is delivered and the package declares no asset dependency.  Editorial formatting only, in addition: an AMS \texttt{amsart} preamble that restates the article's own declarations and macros used by the retained lines, namely the environments \texttt{proposition} on separate counters, together with the standard packages those lines need.  The retained environments print in this document as Proposition 1 in order of appearance, whereas the article prints the same items with the numbering of its own full document.  The complete native source of the article as distributed in the arXiv e-print for this exact version is bundled unchanged as \texttt{sources/180795/source0.tex}.
\begin{align}
\label{eq:gen_lift_fom}
    \dot{\q}&=\p, \nonumber\\
    \dot{\p}&=\D\q -\bar{\kappa} \w_1\odot \w_2,\nonumber \\
    \dot{\w}_1&=\frac{\kappa^2\bar{\kappa}}{2}\w_2 \odot \p, \\
    \dot{\w}_j&=\left(\alpha_j\q +\sum_{i=1}^{k}\alpha_{j,i}\w_i \right) \odot \p,  \qquad \text{for} \qquad j=2, \ldots,k. \nonumber 
    \end{align}

\begin{equation}
\label{eq: lift_energy}
E_{\text{lift}}(\q,\p,\w_1, \cdots,\w_{k})=\frac{1}{2}\p^\top\p - \frac{1}{2}\q^\top \D\q +\frac{1}{\kappa^2} \w_1^\top\w_1. 
\end{equation}

\begin{align} 
\label{eq:learned_lift_rom}
    \dot{\qhat}&=\phat,\nonumber \\
    \dot{\phat}&=\Dhat\qhat + \Hhat_{\p}(\what_1\otimes\what_2) ,\nonumber \\
    \dot{\what}_1&=\Hhat_{\w_1}(\what_2\otimes\phat),\\
        \dot{\what}_j&=\Hhat_{\w_j}(\qhat\otimes\phat) + \sum_{i=1}^k\Hhat_{\w_{j,i}}(\what_i\otimes \phat),  \qquad \text{for} \qquad j=2, \ldots,k. \nonumber 
    \end{align}

\begin{proposition} The structure-preserving Lift \& Learn ROM~\eqref{eq:learned_lift_rom} conserves the perturbed lifted FOM energy 
\begin{equation}
\label{eq:pertrubed_energy}
\widehat{E}_{\mathrm{lift}}(\qhat,\phat,\what_1, \cdots,  \what_{k}):=E_{\mathrm{lift}}(\bPhi\qhat,\bPhi\phat,\V_{1}\what_1, \cdots,\V_{k}\what_{k}) + \Delta E_{\mathrm{lift}}(\qhat)
\end{equation}
with a perturbation of the form $\Delta E_{\mathrm{lift}}(\qhat)=\frac{1}{2}\qhat^\top (\bPhi^\top\D\bPhi-\Dhat)\qhat$ where $\D=\D^\top \in \real^{n \times n}$ is from~\eqref{eq:gen_lift_fom}.
\end{proposition}

\begin{proof}
Given a lifted quadratic FOM of the form~\eqref{eq:gen_lift_fom}, we evaluate the lifted FOM energy function $E_{\text{lift}}$~\eqref{eq: lift_energy} at the lifted state approximation to obtain
\begin{align}
\label{eq:lift_energy}
E_{\text{lift}}(\bPhi\qhat,\bPhi\phat,\V_{1}\what_1, \cdots,\V_{k}\what_{k})&=\frac{1}{2}(\bPhi\phat)^\top(\bPhi\phat) - \frac{1}{2}(\bPhi\qhat)^\top \D(\bPhi\qhat) + \frac{1}{\kappa^2}(\V_{1}\what_1)^\top(\V_{1}\what_1) \nonumber\\
&=\frac{1}{2}\phat^\top\phat - \frac{1}{2}\qhat^\top (\bPhi^\top\D\bPhi )\qhat + \frac{1}{\kappa^2}\what_1^\top\what_1.
\end{align}
Defining $\Delta E_{\mathrm{lift}}(\qhat):=\frac{1}{2}\qhat^\top (\bPhi^\top\D\bPhi-\Dhat)\qhat$ we see that the energy associated with the lifted ROM~\eqref{eq:learned_lift_rom} is--analogous to~\eqref{eq: lift_energy}--given by
\begin{equation*}
\widehat{E}_{\text{lift}}(\qhat,\phat,\what_1, \cdots,  \what_{k})=\frac{1}{2}\phat^\top\phat - \frac{1}{2}\qhat^\top (\bPhi^\top\D\bPhi )\qhat + \frac{1}{\kappa^2}\what_1^\top\what_1 + \frac{1}{2}\qhat^\top (\bPhi^\top\D\bPhi-\Dhat)\qhat=\frac{1}{2}\phat^\top\phat - \frac{1}{2}\qhat^\top \Dhat\qhat +\frac{1}{\kappa^2} \what_1^\top\what_1. 
\end{equation*}
We now compute the time-derivative of $\widehat{E}_{\text{lift}}(\qhat,\phat,\what_1, \cdots,  \what_{k})$
\begin{align*} 
    \frac{\dd}{\dd t}\widehat{E}_{\text{lift}}(\qhat,\phat,\what_1, \cdots,\what_{k})&=\phat^\top\phatdot - (\Dhat\qhat)^\top\qhatdot +\frac{2}{\kappa^2} \what_1^\top\dot{\what}_1 \\
    &=-\bar{\kappa}\phat^\top\bPhi^\top\left( \V_{1}\what_1 \odot  \V_{2}\what_2\right) +  \bar{\kappa} \what_1^\top\V_{1}^\top \left( \V_{2}\what_2 \odot \bPhi\phat\right)\\
    &= 0. 
        \end{align*}
Thus, the quadratic ROM learned via structure-preserving Lift \& Learn conserves $\widehat{E}_{\text{lift}}(\qhat,\phat,\what_1, \cdots,\what_{k})$.
\end{proof}

\end{document}
